\section{The Disputer of This Age}

Few saints are so thoroughly documented, and so persistently contested,
as Robert Bellarmine. He wrote an autobiography---a small, dry,
third-person memoir of 1613 that his own beatification process later
had to defend. His greatest book was, for a season, set down in the
draft of the Roman Index of forbidden books by the very pope to whom
its first volume was dedicated. He was made a cardinal, in the words
his canonization letter attributes to Clement VIII, ``because the
Church of God had not his equal in learning''---and made one, the same
letter insists, \latin{invitum et frustra reluctantem}, unwilling and
resisting in vain. He spent three years as a residential archbishop
and counted them, by every surviving account including his own, among
the best of his life. He signed, with his own hand, the most argued-over
certificate in the history of science---the paper of 26~May 1616 that
Galileo produced in his defense seventeen years later. And he waited
longer for the honors of the altar than almost any comparably famous
candidate in modern history: beatified in 1923, three hundred and two
years after his death; canonized in 1930; declared a Doctor of the
Universal Church in 1931, on the three hundred and tenth anniversary
of the day he died.

This biography asks three questions. First, what do the controlled
sources---Bellarmine's autobiography and letters, the printed volumes
of the \work{Disputationes de controversiis}, the curial and trial
records edited in the national edition of Galileo's works, and the
official acts of 1923--1931---actually establish about the life of
Roberto Francesco Romolo Bellarmino (1542--1621), Jesuit, professor,
consultor, cardinal, archbishop, and spiritual writer? Second, where
the memory is contested---the temporary listing of his own work on
the unpromulgated Sixtine Index; his part in the Roman proceedings
against Giordano Bruno; the warning, the disputed precept, and the
certificate of 1616; the three-century delay of his cause---what does
each surviving document say, and where exactly does document end and
interpretation begin? Third, how did the Church's own judgment of him
move, from the seventeenth-century vetoes of cardinals and courts to
the canonization decretal \work{Lux illa} and the doctoral letter
\work{Providentissimus Deus}, and what do those acts claim---and not
claim---about the man?

The method is that of the whole series: earliest controlled evidence
first; witnesses separated by decades or centuries never silently
merged; the subject's self-presentation read as self-presentation;
legends preserved and labeled rather than suppressed or believed;
official and liturgical reception reported as reception, within the
limits of each act's object. The Galileo matter in particular is
handled at document level: every text quoted there is cited from an
identified witness in Antonio Favaro's national edition, and every
interpretation beyond the texts is named as interpretation. The
terminal appendix states the conventions, the detailed chronology,
and the limits that govern every page.
